question: Explain what the penalty-stage-only structural model is, including any theoretical model of regulator and plant behavior, estimation strategy, how my data is used in the estimation strategy. And where this paper could publish if it answers a penalty-stage-only approach to the question "what are the costs CWSs implicitly place on different violation types, and what does this reveal about the effectiveness of the enforcement regime?"

1. The Theoretical Model
Setting
A CWS operates in each period t facing a latent contamination level driven partly by mining activity and partly by other factors (source water quality, treatment plant aging, etc.). It interacts with a regulator who enforces SDWA standards. The CWS cannot freely choose whether to have contaminated water — the contamination is partly forced on it by upstream mining — but it can choose how to respond to that contamination: fix it (comply), hide it (MR violation), or let it show (MCL violation).

The regulator, after detecting a violation, follows an escalating enforcement process. The CWS plays against this process dynamically.

The CWS's Decision Problem
Each period, conditional on its contamination state s_t = (contamination level, current enforcement status), the CWS chooses an action:


a ∈ { comply, MR violation, MCL violation }
The per-period payoff from each action:

Action	Cost	Mechanism
Comply	-k(s_t)	Pay abatement cost, which rises with contamination level
MR violation	-k_MR	Cheap — just don't file monitoring reports. But triggers audit/inspection with probability π_MR
MCL violation	-k_MCL(s_t)	Contamination exceeds the standard. Inspection probability π_MCL > π_MR; penalty when caught is larger
The CWS's value function:


V(s_t) = max_{a} { u(s_t, a) + β · E[ V(s_{t+1}) | s_t, a ] }
The key economics: MR violations are cheaper than MCL compliance in the short run, but they start the enforcement clock. Once the regulator begins formal enforcement (issues a consent order, mandates treatment upgrades), the CWS faces a large remediation cost. The CWS therefore trades off the current savings from MR violations against the expected present value of getting caught in the penalty machine.

This is the origin of your observed MR substitution: when mining raises a CWS's latent contamination so that MCL compliance becomes expensive, CWSs shift toward MR violations — they stop testing rather than report contamination above the limit — because the expected cost of MR violations is still below the cost of MCL remediation.

The Regulator as a "Regulatory Machine"
Following Duflo et al. exactly, the regulator in the penalty stage is not solving an optimization problem — it is treated as a Markov policy function estimated non-parametrically from the data. The regulator's state is the CWS's violation history (which violation type, how long in violation, prior enforcement actions), and the regulator maps states to action probabilities:


Pr(enforcement_action | violation_history, CWS_characteristics)
Actions available to the regulator: issue informal notice → formal administrative order → penalty → consent order → mandate treatment upgrade → return-to-compliance determination.

This is estimated first, before the CWS's problem, and held fixed in counterfactuals. The philosophical reason to do this: you want to separate what the regulator actually does (descriptive, estimated from data) from what it should do (normative, your counterfactual). The CWS takes the regulator's policy as given when making its compliance decisions.

Information Structure
The regulator does not observe the CWS's latent contamination level directly. It observes:

Whether monitoring reports were filed (MR violation if not)
Contamination readings when tests are submitted (MCL violation if above limit)
Sanitary survey findings (inspections that visit the treatment plant physically)
This creates the CWS's strategic opportunity: by incurring MR violations instead of testing, the CWS prevents the regulator from observing its MCL-level contamination. The CWS knows its own contamination level (it knows what the water looks like); the regulator does not unless it inspects or receives a test result. This information asymmetry is what makes the penalty-stage problem interesting — MR violations are not just cheap, they are informative suppression.

2. Estimation Strategy
The model is estimated in two sequential steps, which is the standard approach for dynamic models with regulatory interactions (Rust 1987; Hotz and Miller 1993; Duflo et al. 2018).

Step 1 — Estimate the Regulatory Machine
Use the administrative SDWA enforcement data to estimate the regulator's Markov transition matrix. For each enforcement state (informal notice, formal order, penalty, consent order, RTC), estimate:


Pr(regulator_action_{t+1} | enforcement_state_t, CWS_characteristics)
This is estimated by maximum likelihood on the enforcement chain sequences in the SDWA data. No structural assumptions are required here — it is pure policy function estimation. The output is a matrix of transition probabilities that summarizes what the regulator actually does in response to each state of the CWS.

This step produces the value of the penalty stage to the CWS: the expected discounted cost V₀(s) of being in enforcement state s, computed by forward-simulating the estimated Markov chain and discounting future enforcement costs. This value function is the key input to the CWS's dynamic problem.

Step 2 — Estimate CWS Compliance Costs Using Conditional Choice Probabilities (CCP)
The CWS's optimal action in each state satisfies:


a*(s_t) = argmax_a { u(s_t, a) + β · E[V(s_{t+1}) | s_t, a] }
Hotz and Miller (1993) showed that the difference in value functions between any two actions can be recovered directly from the observed choice probabilities — you don't need to solve the full Bellman equation. Specifically, the ratio of the probability of choosing MR violation to the probability of complying is a sufficient statistic for the ratio of their costs, conditional on the state.

The structural cost parameters (k_MR, k_MCL, k_comply as a function of contamination level) are recovered by inverting the observed choice probabilities through the model. The key equation:


log[ Pr(MR | s) / Pr(comply | s) ] = [ u(s, MR) - u(s, comply) ] / σ_ε
                                    + β · [ E[V(s')|s, MR] - E[V(s')|s, comply] ] / σ_ε
The left side is observed from data. The second term on the right — the difference in continuation values — is computed from the Step 1 regulatory machine estimates. The only unknowns are the per-period cost differences on the right.

The identification problem: The observed contamination state s_t is endogenous — CWSs in worse source water areas also differ on unobservables that affect their compliance decision. Without an instrument, you cannot tell whether a CWS chooses MR violations because contamination is high (structural) or because it has unobserved characteristics that make it noncompliant generally.

Where the Instrument Enters — the Key Step
The ARP × sulfur instrument enters the estimation here. It provides exogenous variation in the CWS's contamination state — shifts in mining production that change the underlying contamination level the CWS faces, for reasons unrelated to the CWS's own compliance culture or resources.

Formally: instrument the CWS's contamination state s_t with Z_t = post95_t × sulfur_unified_h. This shifts the distribution of latent contamination across CWSs in a way that is:

Exogenous: ARP Phase I timing and coal seam sulfur content are unrelated to the CWS's unobserved compliance characteristics
Relevant: High-sulfur CWSs near mines saw large falls in mining production after 1995, shifting their contamination burden
This gives you two points on the CWS's cost curve: the same CWS before and after the ARP shock faces different contamination levels for exogenous reasons. The shift in its violation choices across these states traces out the cost schedule k(s). A CWS that substitutes more toward compliance when its contamination burden falls exogenously has revealed that its MCL compliance cost fell enough to be worthwhile — which pins down the slope of k(s) relative to k_MR.

In practice this is implemented as an IV/control function approach within the CCP estimator: include the instrument as a first-stage predictor of contamination state, then use the predicted contamination to form the state variable entering the CCP equations.

3. How Your Data Is Used
Data	Role in Estimation
prod_vio_sulfur.parquet (PWSID × year)	Main analysis sample. Violation type outcomes (MCL, MR, share of year in violation) are the observed choice probabilities for the CCP estimator. production_short_tons_coal_unified and sulfur_unified enter as contamination state predictors and instrument components. post95 × sulfur_unified is the instrument Z_t.
SDWA enforcement chain data (ECHO)	Step 1 — estimate the regulatory machine. You need sequences of: violation detected → notice of violation → formal enforcement action → consent order → penalty → RTC determination. Each transition is one observation for the Markov chain MLE.
num_facilities, POPULATION_SERVED_COUNT, PRIMARY_SOURCE_CODE	CWS characteristics X_j that enter the regulatory machine and the cost function — control for scale, source water type (surface water CWSs are more vulnerable to mining runoff).
Sanitary survey frequency (if in SDWA data)	This is the regulator's inspection rate — it enters the regulatory machine as a state variable (CWSs under active sanitary survey scrutiny face different transition probabilities than those that haven't been visited recently).
The moment-matching logic: The model's predicted violation choices (MCL share, MR share, compliance rate) must match the observed violation shares in the data, conditional on the instrumented contamination state. The instrument allows this matching to be done at exogenously-shifted contamination levels rather than endogenously-determined ones.

One practical data challenge: The SDWA enforcement chain needs to be dense enough to estimate Markov transition probabilities with reasonable precision. If the enforcement data is sparse (many CWSs with few enforcement interactions), you may need to pool across CWS types or use a simpler Markov structure (fewer states). You noted from a prior session that the enforcement chain data (sanitary surveys → formal enforcement → RTC) is usable for H2 and H3 outcomes — which suggests there is enough data for the regulatory machine estimation, but you would want to verify the number of enforcement chain sequences before committing.

4. What the Model Reveals About Enforcement Effectiveness
The structural estimates directly answer:

Implicit cost ratios. The ratio k_MR / k_MCL recovered from the CCP estimates tells you how much cheaper CWSs find MR violations relative to MCL remediation. If this ratio is very low (MR violations are essentially free to incur), it reveals that the current enforcement regime creates almost no deterrence against monitoring suppression — the penalty machine for MR violations does not generate enough expected cost to induce reporting.

Penalty machine effectiveness. From Step 1 you observe directly whether the regulator actually escalates enforcement after MR violations — whether transition probabilities from MR violation to formal enforcement are high or low. If the regulator rarely escalates, the expected cost of MR violations is structurally low, explaining the substitution.

Counterfactuals:

Increase the MR penalty transition probability (simulate a more aggressive enforcement regime) → how much does the MR-MCL substitution fall?
Increase sanitary survey frequency for high-contamination CWSs → how much does MCL compliance improve?
Compare actual enforcement outcomes to the efficient enforcement benchmark — what fraction of the MCL compliance gap is attributable to the low cost of MR violations rather than the intrinsic cost of MCL remediation?
5. Publication Venues
The combination of quasi-experimental identification + structural model for policy counterfactuals is the dominant methodological paradigm at top journals right now. Where the paper lands depends on how much the structural model adds to what you can say with reduced form alone.

If the structural model is the main contribution:

Venue	Fit	Reason
JPE	Strong	Structural regulatory models + clean quasi-experimental identification is JPE's sweet spot (Ryan 2012 on Clean Air Act is the archetype)
AER	Strong	If the counterfactuals are policy-relevant and the identification is airtight
Review of Economic Studies	Strong	Structural dynamic models are well-received; RES has published environmental regulatory models
AEJ: Applied	Very strong	More receptive to applied structural work with a quasi-experimental identification step; the combination is exactly what AEJ: Applied has been publishing
If reduced form is the main contribution and structural is a section:

Venue	Fit	Reason
AEJ: Policy	Very strong	Policy-relevant enforcement question, SDWA is understudied, clean instrument
Journal of Public Economics	Strong	Public regulation + enforcement + environmental externalities
JEEM	Strong	Natural home for environmental enforcement papers; structural enforcement models have published there
RAND Journal	Moderate	IO venue; structural regulatory models fit, but RAND prefers industrial organization settings more narrowly
The honest assessment of ceiling:

A penalty-stage-only structural model without the targeting stage is a significant step below the Duflo et al. full paper. The reason Duflo et al. is in Econometrica is the joint identification of the targeting and penalty stages — the experiment lets them separately measure the value of the regulator's private information, which is the novel theoretical result. A penalty-stage-only model recovers cost parameters but does not identify the regulator's information problem, so the theoretical contribution is smaller.

However: your setting has one advantage Duflo et al. lacks — you have a geophysical instrument for contamination that is completely outside the regulatory-CWS interaction. Their experiment was necessary because they had no instrument for the pollution state. Your ARP × sulfur instrument means you achieve clean identification of compliance costs without needing the experimental variation. That is a meaningful methodological contribution in its own right, and it could be framed as: "we show that quasi-experimental variation in the regulated entity's pollution burden — rather than randomized inspections — is sufficient to identify structural compliance costs, and we apply this to the understudied SDWA enforcement context."

That framing puts AEJ: Applied or Journal of Public Economics as the realistic ceiling for a standalone structural paper, with a genuine shot at AER if the counterfactual policy results are striking (e.g., you show that the current enforcement regime recovers only X% of the potential compliance gains relative to a penalty-optimal regime).

question: I would like comparative statics for my theoretical model. How can I get comparative statics from the model adaptation of Z:\ek559\mining_wq\.claude\skills\Econometrica - 2018 - Duflo - The Value of Regulatory Discretion  Estimates From Environmental Inspections in India.pdf. If I cannot get comparative statics from Z:\ek559\mining_wq\.claude\skills\Econometrica - 2018 - Duflo - The Value of Regulatory Discretion  Estimates From Environmental Inspections in India.pdf can the theoretical model from Z:\ek559\mining_wq\.claude\skills\Kang and Silveria 2021.pdf be applied to the suggested plan to adapt Z:\ek559\mining_wq\.claude\skills\Econometrica - 2018 - Duflo - The Value of Regulatory Discretion  Estimates From Environmental Inspections in India.pdf to get comparative statics or is it too different from the structural model I would actually be estimating to be useful?

Question 1: Can you get comparative statics from the Duflo et al. penalty-stage-only model?
No — not analytically. Only numerically.

The reason is structural. The Duflo et al. penalty stage is a dynamic discrete choice (DDC) model in which the CWS's value function V(s) is defined implicitly by the Bellman equation:


V(s) = max_a { u(s,a) + β · E[V(s')|s,a] }
There is no closed-form expression for V(s) as a function of the underlying parameters. If you try to differentiate this equation with respect to a parameter ψ (say, the penalty severity), you get:


∂V/∂ψ = (∂u/∂ψ) + β·(∂E[V']/∂a*)·(∂a*/∂ψ) + β·(∂E[V']/∂ψ)
The term ∂a*/∂ψ — how the optimal action changes with the parameter — depends on the curvature of V, which itself depends on ∂V/∂ψ. The system is circular and has no closed-form resolution without specific parametric assumptions strong enough to make the model trivial.

What you get instead from the Duflo framework is numerical counterfactuals: estimate the model, change a parameter, re-solve computationally, compare equilibria. These tell you how much compliance changes under alternative penalty regimes, but they do not generate theorems about why it changes in a direction you can state and prove before estimation. They are quantitative policy simulations, not comparative statics.

This is not a flaw of the Duflo model — it is a feature. The model is rich enough to capture multi-period enforcement chains precisely because it doesn't collapse to a tractable closed form. But it means the Duflo adaptation cannot serve as the source of your comparative statics propositions.

Question 2: Can Kang & Silveria (K&S) provide comparative statics?
Yes — K&S yields explicit, provable comparative statics propositions.

The reason is that K&S uses the Mookherjee-Png (1994) model, which is a static mechanism design / screening framework with a well-defined first-order condition for the facility's problem.

The K&S model structure (adapted to your setting)
Each CWS has a private compliance cost type θ — the cost of treating contaminated water to meet MCL standards. Higher θ means MCL compliance is more expensive (e.g., the CWS's source water is more contaminated, partly due to upstream mining).

The CWS chooses a negligence level a ∈ [0, ā], where higher a means less monitoring effort — operationally, this maps to the choice to suppress monitoring (MR violations) rather than submit readings that would reveal MCL-level contamination.

The number of violations K follows a Poisson distribution with mean a. The CWS faces an expected penalty:


e(a) = exp(-a) · Σ_k ε(k)/k! · aᵏ
where ε(k) is the penalty schedule the regulator sets as a function of observed violations k.

The CWS's payoff from choosing negligence level a, given cost type θ, is:


-θ[b(ā) - b(a)] - e(a)
The first term is the compliance cost avoided by choosing a over full compliance (θ scales the private benefit of negligence); the second is the expected penalty. The CWS maximizes this, giving the first-order condition (K&S equation 3):


θ · b'(a(θ)) = e'(a(θ))
This single equation in one unknown (a), as a function of two parameters (θ and the penalty schedule e), is the source of all comparative statics. It can be implicitly differentiated.

The three comparative statics propositions
Proposition 1 — Type sorting (the selection mechanism).

Differentiate the FOC with respect to θ:


b'(a) + θ·b''(a)·(∂a/∂θ) = e''(a)·(∂a/∂θ)
∂a/∂θ = b'(a) / [e''(a) - θ·b''(a)]
Under the second-order condition (e''(a) - θ·b''(a) > 0) and since b'(a) > 0 (more negligence saves more costs):

∂a/∂θ > 0: Higher-compliance-cost CWSs choose more negligence — they incur more MR violations and fewer MCL-compliance-driven remediation actions in equilibrium. This is the theoretical foundation for why the MR-MCL substitution pattern is a CWS-level response to heterogeneous compliance burdens, not noise.

Proposition 2 — Enforcement deterrence.

Suppose the penalty schedule shifts upward by λ·m(a), representing stricter enforcement. Differentiate the FOC with respect to λ:


0 = [e''(a) + λ·m''(a)]·(∂a/∂λ) + m'(a)
∂a/∂λ = -m'(a) / [e''(a) + λ·m''(a)] < 0
∂a/∂λ < 0: A uniform increase in enforcement stringency reduces negligence — strictly more compliance — for every type θ. This is the deterrence comparative static: steeper penalty schedules reduce the MR violation equilibrium.

Proposition 3 — The mining externality as a type shifter.

In your setting, θ is not fixed — it is partly determined by upstream mining production. Higher mine production increases the contamination burden of downstream CWSs, making MCL compliance more expensive. Write θ = θ(m) where m is mine production and ∂θ/∂m > 0.

By the chain rule on Proposition 1:


∂a*/∂m = (∂a*/∂θ) · (∂θ/∂m) > 0
More mining → higher effective compliance cost type → higher equilibrium negligence → more MR violations. This is the theoretical comparative static that your 2SLS reduced form is testing empirically: the ARP shock reduces m, shifting θ(m) down for high-sulfur CWSs, which moves a* downward (less negligence, more compliance), producing the reduced-form effect you already estimate.

Proposition 4 — Optimal penalty convexity (requires full K&S regulator optimality).

The regulator's problem (K&S equation 4) minimizes the sum of expected compliance costs, environmental harms from violations, and enforcement costs, by choosing ε(k). The optimality conditions imply that the optimal penalty schedule is strictly convex in the number of violations: each additional violation triggers a larger marginal penalty increment. This is the K&S main result — it justifies why penalty schedules that escalate super-linearly in violation counts are efficient screening devices.

Note: Proposition 4 requires imposing that the regulator is optimally designing the penalty schedule. Your estimation holds the regulatory machine fixed (exogenous), so you cannot use Proposition 4 as a testable implication of your estimated model. It remains a normative benchmark for counterfactuals.

Question 3: Is K&S too different from the Duflo estimation strategy to be useful?
No — it is not too different, but the mismatch requires explicit acknowledgment. Propositions 1–3 carry over; Proposition 4 does not.

Here is the precise accounting of what is and is not compatible:

What carries over
The type sorting and enforcement deterrence logic (Propositions 1–3) are driven entirely by the CWS's own optimization — the FOC θ·b'(a) = e'(a). This FOC does not depend on the regulator's optimality. It holds as long as the CWS is maximizing its own payoff given any enforcement function e(a), whether or not the regulator designed e(a) optimally.

In the dynamic Duflo model, the CWS's first-order condition is a dynamic analogue — the CWS equates the marginal benefit of negligence against the marginal increase in the continuation value of being caught. The direction of comparative statics is the same: more negligence is more attractive when θ is high (Proposition 1) and less attractive when enforcement is more severe (Proposition 2). The dynamic model adds the feature that the comparison is over continuation values rather than static penalties, but it does not reverse any sign.

The mining externality (Proposition 3) is purely an exogenous shift in the distribution of θ across CWSs — it does not depend on any feature of the dynamic model at all.

What does not carry over
Proposition 4 (optimal convexity) requires that the regulator is optimally solving the mechanism design problem in K&S equation 4. In the Duflo adaptation, you explicitly do not impose this — the regulatory machine is estimated as a descriptive Markov process, not derived from a regulator optimization. If you imposed K&S's regulator optimality on your model, you would be assuming the answer to one of the questions your paper is trying to ask (whether the enforcement regime is well-designed). Proposition 4 is useful as a normative benchmark for counterfactuals ("how much better would compliance be if the penalty schedule were the K&S optimum?") but cannot be presented as a testable prediction of your estimated model.

The standard practice for handling this mismatch
This gap between a tractable static theoretical model and a richer dynamic empirical model is entirely normal and handled uniformly in the literature:

Ryan (2012) in JPE derives comparative statics from a static entry game for intuition, then estimates a dynamic model for counterfactuals.
K&S itself is a static model, yet the enforcement processes it describes are clearly multi-period.
Duflo et al. explicitly notes they abstract from dynamics in certain parts of their theoretical discussion.
The standard framing in the theory section: "We present a static model of regulator–CWS interaction to derive comparative statics predictions. The empirical model extends this to a dynamic setting to account for the multi-period enforcement chain observed in the data. Propositions 1–3 hold qualitatively in both frameworks; Proposition 4 serves as a normative benchmark."

The combined architecture
The cleanest paper structure uses both models in sequence:

Paper Section	Framework	What It Provides
Theory	K&S / Mookherjee-Png (static)	Propositions 1–3 with proofs; Proposition 4 as normative benchmark
Reduced form	2SLS (ARP × sulfur)	Tests Proposition 3 empirically — mining increases violations
Structural estimation	Duflo penalty stage (dynamic)	Recovers k_MR / k_MCL cost ratio; estimates regulatory machine
Counterfactuals	Re-solve Duflo model	Quantifies Propositions 2 and 4 numerically
This is a coherent, internally consistent architecture. The theory section gives you propositions you can prove before looking at the data. The reduced form tests one of them directly. The structural model recovers the primitive cost parameters that underlie all of them. The counterfactuals simulate the quantitative magnitude of the deterrence and optimality propositions. None of the sections contradict each other — they are complementary tools for answering the same question at different levels of rigor and richness.